% !TEX root = ../3.GSM.tex
\section{Experimental Method}

\subsection{Mill}
\begin{itemize}
    \item Turn on the machine without load to measure the no-load power.
    \item Give 200\,g of rice to the mill.
    \item Turn on the input screw and press the meter to measure the maximum load current intensity.
    \item When the current returns to the no-load value, press the meter to determine the grinding time.
\end{itemize}

\subsection{Sieve}
\begin{itemize}
    \item Determine sieve efficiency at the 0.2\,mm sieve size: sieve 80\,g of ground rice 5 times, each time for 5 minutes, and weigh the rice passing through the sieve each time.
    \item Determine the size distribution of the material after grinding: sieve 80\,g of ground rice through various sieves of different sizes for 20 minutes and weigh the accumulated rice in each sieve.
\end{itemize}

\subsection{Mix}
\begin{itemize}
    \item Put 1.5\,kg of green beans and 3\,kg of soya beans into the machine.
    \item Turn on the machine.
    \item Stop the machine at 6 time intervals (5\,s, 15\,s, 30\,s, 60\,s, 120\,s, 300\,s) and take 8 samples at each time according to the sampling chart. Count each particle in the sample.
\end{itemize}

The 8 sample positions are arranged in the mixing vessel as follows:

\begin{center}
\begin{tabular}{ccc}
    1 & 2 & 3 \\
      & 4 & 5 \\
    6 & 7 & 8 \\
\end{tabular}
\end{center}
